\chapter[Limits of Certification]{The Limits of Certification and the Future of Mathematical Understanding}
\label{ch:conclusion}

The monograph began with a proof that compiled and an argument that it did not follow. It ends by asking what it means to know that a mathematical problem has been solved when the proof is formally valid, the problem is expressed in natural language, and the relation between them may not be completely decidable. The question has two familiar answers, and the monograph rejects both.

\section{Two answers that fail}

The first answer holds that formal verification settles every relevant epistemic question. On this view a proof that the kernel accepts is a proof, and anything further is superstition about trusted bases. The answer fails because the kernel's jurisdiction is a relation between a term and a type in a stated environment, and the questions of Chapter~\ref{ch:nondischarge} lie outside it. Kernel acceptance is blind to the specification, to the argument, and to the interpretation, and no refinement of the kernel can change that, since the blindness is a property of what the judgment is a judgment of. The hidden-existence family shows that the questions it leaves open are not merely unattended but, on unrestricted classes, hard at every level of the arithmetical hierarchy for well-definedness, and hard at the second level for correspondence under the stated convention.

The second answer holds that because semantic uncertainty cannot be eliminated, machine-generated and machine-verified mathematics is worthless, or at least cannot be trusted. This answer fails for a reason developed in Chapter~\ref{ch:futility}. The existence of undecidable instances limits a single total procedure and leaves a large space of sound, incomplete, and restricted procedures. A proof that compiles is genuine evidence about the formal statement. The two Navier-Stokes cases do not show that the Lean development is worthless, and they do not show that the natural-language argument is wrong. They show that the relation between the two is unsettled and can be described precisely. To conclude from uncertainty that nothing has been learned is as unwarranted as to conclude from acceptance that everything has.

\section{Bounded certification}

The alternative is a theory of bounded certification. Its requirements are three, and each has a counterpart in the formal development. A system should state exactly what it establishes, which is the discharged obligations with their currently warranted certificates, the authorities, the interpretation, and the dependencies. It should preserve what it cannot establish, which is the entailment condition of Chapter~\ref{ch:preservation}, under which unresolved obligations are carried through transformations without being relabeled as settled. And it should expose the conditions under which further evidence could change a result's status, which is the dependency graph with its defeaters, so that for each discharged obligation the system can list what would revoke the discharge, and for each open one what would settle it.

The three requirements share a feature. None of them eliminates uncertainty. Each makes the uncertainty a structured object that can be inspected, transported, and reduced. The residue of a verification is a precise list of what remains to be done, and a corpus with a small, benign, well-described residue is more trustworthy than a corpus with no recorded residue, because the second kind cannot distinguish a settled question from an unasked one.

\section{Preservation is not discharge}

The distinction that the monograph regards as its contribution is that between preserving an obligation and discharging it. A pipeline can rigorously carry an unresolved semantic obligation through a sequence of translations and revisions, certifying at each step that nothing has been lost, and arrive at the end with the obligation exactly as open as it began. The preservation is a mathematical fact with a proof, Theorem~\ref{thm:compose}, conditional on the soundness of its certificates. The openness is a fact about the evidence. Preservation does not imply discharge, and discharge of one obligation does not secure preservation of another, and Proposition~\ref{prop:independence} shows both can occur. An obligation recorded in the discharge set of a transition is preserved at that transition, so the independence is asymmetric at a single recorded discharge. A culture of verification that treats preservation as discharge claims too much, and one that treats the absence of discharge as the absence of preservation discards what has been secured.

The distinction supports a modest and a more ambitious reading. The modest reading is a discipline of reporting: say what was checked, say what was carried, and say what remains. The ambitious reading is a research program in which the certificates that justify entailment and discharge are themselves objects of study, with the questions listed at the end of Chapter~\ref{ch:preservation} concerning the certification of entailments between natural-language and formal contents, the justification of refinements, and the changing of interpretations along a chain. These remain open.

\section{Understanding}

The title of the chapter refers to the future of mathematical understanding, and the monograph has said little about understanding directly. It has treated the judgment of correspondence as an act performed by identified authorities, recorded as certificates, and defeasible. It has not treated understanding as a verifier that can be inserted at one point of a pipeline and discharge the semantic residue there. Chapter~\ref{ch:reward} argued that understanding at the point where objectives are set is necessary and not sufficient, because the interpretation of novel problems is constructed after the objective is reviewed. The more general lesson is that understanding is not located at one stage. It is exercised wherever a correspondence is asserted, and a system that makes those assertions visible and attributable distributes the burden of understanding where it belongs and makes it possible to see where it has not been exercised.

It would be a mistake to read the monograph as a pessimistic account of formal methods. Its conclusions about what verification does not establish are conclusions about the limits of a jurisdiction, and a jurisdiction whose limits are known is more useful than one whose limits are assumed to coincide with those of the mathematical problem. The work closes by distinguishing the elimination of uncertainty, which the results of Part~III show to be unavailable on unrestricted classes, from its responsible representation, which is available and which this monograph has tried to specify. Verification becomes more trustworthy when the system no longer mistakes the limits of its own jurisdiction for the limits of the mathematical problem.
